\section{GPT-3: meta-learning y curvas de aprendizaje few-shot}

GPT-3 va más allá y coloca examples dentro del context, de modo que el modelo ajusta su salida según la demonstration local sin actualizar parámetros. La clase lo llama language-model meta-learning: el outer loop es el preentrenamiento y el inner-like behavior ocurre en el context. Esta formulación es una analogía funcional; no hace falta suponer que el modelo realmente ejecuta descenso de gradiente.

\subsection{Meta-learning interface}

Esta sección define primero la interfaz de la tarea en tiempo de inferencia. \term{few-shot} se refiere a proporcionar dentro del prompt unos pocos input-output examples; \term{in-context learning} se refiere a que el modelo cambia su comportamiento únicamente con base en ese contexto. Para los parámetros $\theta$, la predicción sigue siendo $p_\theta(y\mid x,\text{demonstrations})$, y $\theta$ no cambia.

\lecturefigure{slide-18-gpt3-metalearning.jpg}{GPT-3 coloca la task description y las demonstrations en el mismo context.}{Video oficial 18:15; Brown et al. 2020.}

\begin{knowledgebox}{Lectura de la figura: la estructura de dos niveles entre entrenamiento e inferencia}
El preentrenamiento actualiza los parámetros sobre una gran variedad de formatos de tarea; el inference prompt aporta una definición temporal de la tarea. El modelo puede producir el behavior mediante pattern completion, retrieval-like matching o un learned algorithm. Sea cual sea el mecanismo interno, en la práctica de ingeniería hay que probar el demonstration order, las label words, el format y el context length.
\end{knowledgebox}

\subsection{Few-shot arithmetic}

Esta sección usa la arithmetic curve para examinar la interacción entre scale y demonstrations. La figura compara zero/one/few-shot con el model size; algunas tareas, según el número de dígitos, mejoran notablemente al aumentar la escala, pero las curvas difieren según la dificultad, lo que indica que la capacidad no es un scalar único.

\lecturefigure{slide-19-few-shot-arithmetic.jpg}{Las curvas de arithmetic de GPT-3 muestran la interacción entre scale, shots y dificultad de la tarea.}{Video oficial 19:23.}

\begin{knowledgebox}{Lectura de la figura: primero, qué representa cada curva}
El eje horizontal suele ser el model size y el vertical la accuracy; el color o el tipo de línea distingue el número de dígitos o el shot setting. Lo importante no es que «más grande siempre es mejor», sino que algunas tareas tienen una pendiente más pronunciada y otras se quedan cerca del azar. Acertar en aritmética también puede depender de la tokenización y del formato de los datos de entrenamiento, por lo que no equivale directamente a razonamiento matemático general.
\end{knowledgebox}

\subsection{Word unscrambling}

Unscrambling pide al modelo recuperar una palabra a partir de caracteres o fragmentos de palabra. Se acerca más a un algorithmic mapping que la language completion habitual, y permite observar si las demonstrations llevan al modelo a inferir la regla.

\lecturefigure{slide-20-word-unscrambling.jpg}{Word unscrambling usa few-shot examples para probar la inducción de reglas.}{Video oficial 20:14.}

\begin{knowledgebox}{Lectura de la figura: atención a la task format sensitivity}
Que la curva suba indica que los modelos más grandes aprovechan mejor los context examples, pero la segmentación por caracteres, las mayúsculas y minúsculas, los espacios y el formato de la respuesta modifican la dificultad. Si basta con cambiar el label token para que el desempeño caiga notablemente, lo que el modelo aprendió puede ser un pattern frágil y no un algorithm estable.
\end{knowledgebox}

\subsection{General few-shot learning}

A continuación, el desempeño en múltiples tareas y la scale se reúnen en una misma figura de síntesis. El Average puede mostrar la tendencia global, pero también oculta las diferencias entre tareas, la transferencia negativa y la benchmark saturation; por eso hay que volver a la per-task evidence.

\lecturefigure{slide-21-general-few-shot.jpg}{Las curvas multitarea de GPT-3 colocan zero/one/few-shot y la tendencia con la escala en los mismos ejes.}{Video oficial 20:36.}

\begin{knowledgebox}{Lectura de la figura: debajo del promedio está la distribución de tareas}
Compare la separación entre zero-shot, one-shot y few-shot, y luego observe si crece con la escala. Si la mejora few-shot proviene solo de unas cuantas tareas, el promedio engaña; además, el benchmark puede estar inflado por contamination del preentrenamiento. Una conclusión confiable requiere un per-task breakdown, confidence y held-out data governance.
\end{knowledgebox}

\teachervoice{El ponente describe el cambio central de GPT-3 como «language model metalearning». La nota de clase se centra en la interface: la tarea se coloca en el context en lugar de reentrenar el modelo para cada tarea. Esta interfaz cobró gran importancia después, pero las curvas de entonces ya mostraban que no funciona de manera uniforme entre tareas y formatos.}

\subsection{Resumen de la sección}

GPT-3 convierte las demonstrations en entrada en tiempo de inferencia y así da forma a la in-context interface. La scale mejoró la capacidad de aprovechar el contexto, pero no eliminó la sensibilidad al formato, las diferencias entre tareas ni el riesgo de contamination.
